\begin{table}[H]
\centering
\small
\caption{BVRP e variação da IV pelo sinal do retorno do dia.}
\label{tab:retnl-sinal}
\begin{tabular}{lrrrr}
\toprule
Variável & $r_t > 0$ & $r_t < 0$ & Diferença & $p$ \\
\midrule
$\text{BVRP}_t$ (p.p.) & $-$9{,}00 & $-$8{,}63 & 0{,}37 & 0{,}564 \\
$\Delta_1\text{IV}_t$ (p.p.) & $-$0{,}46 & 0{,}41 & 0{,}88 & $<$0{,}001 \\
\addlinespace
\multicolumn{5}{l}{$\Delta_1\text{IV}_t$ sobre $r_t^+$ e $r_t^-$: $\hat b^+ = 0{,}005$ ($p = 0{,}958$); $\hat b^- = -0{,}772$ ($p < 0{,}001$)} \\
\multicolumn{5}{l}{Simetria, $b^+ + b^- = 0$: $p < 0{,}001$} \\
\bottomrule
\end{tabular}
\par\smallskip
\parbox{0.95\linewidth}{\footnotesize\textit{Nota}: Amostra de modelagem; 877 dias com $r_t > 0$ e 899 com $r_t < 0$. Diferença: dias de queda menos dias de alta, pela regressão numa indicadora de $r_t < 0$, com HAC de 31 defasagens para o BVRP (alvos de 30 dias sobrepostos) e de 2 para $\Delta_1\text{IV}_t$. Duas últimas linhas: $\Delta_1\text{IV}_t = a + b^+ r_t^+ + b^- r_t^- + u_t$, com $r_t^+ = \max(r_t, 0)$ e $r_t^- = \min(r_t, 0)$ em \%, HAC de 2 defasagens; a simetria é $b^+ = -b^-$.}
\end{table}
